\section{Kompilierungsprobleme}

Erste Hilfe, wenn etwas nicht kompiliert (also wenn der Code nicht korrekt in eine
pdf umgewandelt wird):\footnote{Folgende Auflistung ist übernommen aus dem Script von Stefan Macke.}

\begin{itemize}
    \item Immer die erste Fehlermeldung lesen, alle weiteren sind meist Folgefehler.
    \item Fehlende schließende Klammer und fehlendes end zur Umgebung sind die häufigsten Ursachen.
    \item Verzeichnisse fehlen oder zeigen Fragezeichen, wenn zu selten kompiliert wurde. Zweimal laufen lassen, bei Literatur zusätzlich biber oder bibtex dazwischen.
    \item Hilfsdateien mit den Endungen aus, toc und log lsöchen, wenn ein Lauf sich hartnäckig falsch verhält.
\end{itemize}